\section*{A003 endpoint interfaces (bits throughout)}
Let $S\sim P_f$, $P_f(A)=\widehat f(A)^2$, $\alpha_i=\Pr(i\in S)$, and
\[
 E_i=H(1_{\{i\in S\}}\mid S\setminus\{i\}),\qquad
 \Delta_iH=H(f)-\frac{H(f_{i,+})+H(f_{i,-})}{2}.
\]
Only coordinates with $\alpha_i>0$ are relevant.

\paragraph{Exact split identities.}
Writing $a_T=\widehat f(T)$ and $b_T=\widehat f(T\cup\{i\})$, the child laws are
$r_\pm(T)=(a_T\pm b_T)^2$ and their midpoint is $q(T)=a_T^2+b_T^2$. Hence
\[
 \Delta_iH=E_i+\operatorname{JS}(r_+,r_-),\qquad
 I(f)-\frac{I(f_{i,+})+I(f_{i,-})}{2}=\alpha_i,
\]
and
\[
 0\le\operatorname{JS}(r_+,r_-)
 \le\frac2{\ln2}\sum_T\min(a_T^2,b_T^2)
 \le\frac2{\ln2}\min(\alpha_i,1-\alpha_i).
\]
The pointwise scalar inequality is
\[
 \frac{(1+t)\log_2(1+t)+(1-t)\log_2(1-t)}2
 \le\frac{1-\sqrt{1-t^2}}{\ln2}\quad(|t|\le1).
\]

\paragraph{R: relevant-coordinate step criterion.}
\[
 \mathsf{STEP}(K):\quad
 \forall n\ \forall f\text{ nonconstant}\ \exists i\ (\alpha_i>0):
 \Delta_iH(f)\le K\alpha_i.
\]
Strong induction gives $\mathsf{STEP}(K)\Rightarrow\mathsf{FEI}(K)$.
Moreover $B(f)=\sum_iE_i\le H(f)$, so
\[
 \mathsf{FEI}(C)\Rightarrow\mathsf{STEP}(C+2/\ln2).
\]
Thus $C_{\rm FEI}\le C_{\rm step}\le C_{\rm FEI}+2/\ln2$.

\paragraph{Signed Booleanity/downward criterion.}
For signed coefficients $a_A$, require the complete system
\[
 \sum_Aa_A^2=1,\qquad \sum_Aa_Aa_{A\triangle U}=0
 \quad(\varnothing\ne U\subseteq[n]).
\]
Set $p_A=a_A^2$ and
\[
 L=\sum_Ap_A\sum_{i\in A}\log_2\left(1+\frac{p_{A\setminus\{i\}}}{p_A}\right).
\]
Pairing $T,T\cup\{i\}$ proves
\[
 L_i\le E_i\le L_i+\alpha_i/\ln2,qquad L\le B\le L+I/\ln2.
\]
Therefore
\[
 \mathsf{DOWN}(K):\ \forall a\text{ satisfying full signed Booleanity},\ L\le KI
 \quad\Longrightarrow\quad \mathsf{FEI}(K+3/\ln2),
\]
while $\mathsf{FEI}(C)\Rightarrow\mathsf{DOWN}(C)$. Hence
$C_L\le C_{\rm FEI}\le C_L+3/\ln2$.

\paragraph{J: complete conditioned posterior residual.}
For admissible ordered gates, fractional weights, and laws $\mathcal P_A$ on
distinct actual finite partition-forest runs producing $A$, define
\[
 \mathcal J_f=
 v(\Lambda+\log_2Z)+\sum_{A\ne\varnothing}P_f(A)
 \left[\mathbb E_{r\sim\mathcal P_A}B_\lambda(A;r)-H(\mathcal P_A)\right],
\]
where $\Lambda=\sum_l\log_2(1+\lambda_l\pi_l)$ and $Z=1-2^{-\Lambda}$.
The recovered theorem is
\[
 H(f)\le h_2(v)+K_{\rm ref}I(f)+\mathcal J_f,qquad
 K_{\rm ref}=3+2\log_2 3+\log_2(6e).
\]
Consequently
\[
 \mathsf{WF}(K):\quad \forall n\ \forall f\ \exists
 (\text{admissible data})\quad \mathcal J_f\le KI(f)
 \Longrightarrow \mathsf{FEI}(2+K_{\rm ref}+K).
\]
No converse from FEI to $\mathsf{WF}$ is known.

\paragraph{Implication graph.}
\[
 \begin{array}{ccccc}
 \mathsf{WF}(K)&\Longrightarrow&\mathsf{FEI}(2+K_{\rm ref}+K)\\[2mm]
 \mathsf{FEI}(C)&\Longrightarrow&\mathsf{STEP}(C+2/\ln2)&\Longrightarrow&
 \mathsf{FEI}(C+2/\ln2)\\[2mm]
 \mathsf{FEI}(C)&\Longrightarrow&\mathsf{DOWN}(C)&\Longrightarrow&
 \mathsf{FEI}(C+3/\ln2).
 \end{array}
\]
The first arrows in the last two rows use $B\le H$ and $L\le B$; the final
arrows use physical restriction induction. The three unpaid leaves are the
dimension-independent bounds $\mathsf{WF}(K)$, $\mathsf{STEP}(K)$, and
$\mathsf{DOWN}(K)$.

\paragraph{Arbitrary-seed admission.}
For $F(x,y)=f(x)g(y)$ on disjoint variables, $H,I,L$ are additive and a
coordinate's $(\Delta_iH,\alpha_i)$ pair is unchanged from its factor.
Therefore tensor powers preserve the obstructing ratios. Each target also
quantifies over every genuine Boolean seed, including sparse defects,
addressing functions, near tensors, and prime gates.
